\section{Spectral Theory in Hilbert Spaces}\label{sec:spectral_theory_in_hilbert_spaces}
In this section we introduce the basic elements of spectral theory for operators on Hilbert spaces and connect them to the description of quantum observables and measurements. Throughout, $\mathcal{H}$ denotes a complex Hilbert space with inner product $\langle\cdot,\cdot\rangle$ conjugate-linear in the first argument and linear in the second. We write $\braket{\psi|\phi}:=\langle\psi,\phi\rangle$ in Dirac notation. The associated norm is
\begin{equation}
\lvert\ket{\psi}\rvert := \sqrt{\braket{\psi|\psi}},
\end{equation}
and, for brevity, we will often write $\lvert\psi\rvert$ instead of $\lvert\ket{\psi}\rvert$.
\subsection{Self-adjoint Operators and Observables}\label{sec:self_adjoint_operators_and_observables}
In quantum mechanics, observables are represented by self-adjoint operators on a Hilbert space.\supercite{1378551877,1423848155} Spectral theory describes how such operators can be decomposed into their eigenvalue parts, both discrete and continuous, and how this decomposition encodes measurement statistics and time evolution.
\begin{definition}[Self-Adjoint Operator]
A linear operator $A$ on $\mathcal{H}$ is called self-adjoint if
\begin{equation}
\braket{\phi|A\psi}=\braket{A\phi|\psi}\quad\text{for all }\ket{\phi},\ket{\psi}\in\mathcal{H}
\end{equation}
and the domain of $A$ coincides with the domain of its adjoint.
\end{definition}
In physics terminology, self-adjoint operators are often called Hermitian. In nonrelativistic quantum mechanics, each observable is modeled by a self-adjoint operator $A$ on $\mathcal{H}$ and its spectral properties determine the possible measurement outcomes.

Since a lot of operators in quantum mechanics require unitarity, we also define it here.
\begin{definition}[Unitary Operator]
A linear operator $U$ on a Hilbert space $\mathcal{H}$ is called unitary if
\begin{equation}
U^\dagger U=U U^\dagger=I,
\end{equation}
where $U^\dagger$ is the adjoint of $U$ and $I$ is the identity on $\mathcal{H}$.
\end{definition}
Unitary operators preserve the inner product and hence the norm
\begin{equation}
\braket{U\psi\vert U\phi}=\braket{\psi\vert\phi}\quad\forall\psi,\phi\in\mathcal{H},
\end{equation}
so that $\lvert U\psi\rvert=\lvert\psi\rvert$. In quantum mechanics they represent symmetry
transformations and the time evolution of isolated systems. If $\mathcal{H}$ is finite-dimensional, the spectral structure is simple.\supercite{1378551877}
\begin{theorem}[Spectral Theorem]
Let $\mathcal{H}$ be a finite-dimensional Hilbert space and $A$ self-adjoint. Then there exists an orthonormal basis $\{\ket{n}\}$ of $\mathcal{H}$ such that
\begin{equation}
A\ket{n}=\lambda_n\ket{n},\qquad\lambda_n\in\mathbb{R},
\end{equation}
and
\begin{equation}
A=\sum_n\lambda_n\ket{n}\bra{n}.
\end{equation}
\end{theorem}
The real numbers $\lambda_n$ are the eigenvalues of $A$ and the kets $\ket{n}$ are the corresponding eigenstates. If the system is in a normalized state $\ket{\psi}$, the probability of obtaining the result $\lambda_n$ when measuring $A$ is
\begin{equation}
\mathbb P_A^\psi(\lambda_n)=\lvert\braket{n\vert\psi}\rvert^2,
\end{equation}
and the expectation value of $A$ is
\begin{equation}
\langle A\rangle_\psi=\braket{\psi\vert A\vert\psi}=\sum_n\lambda_n\lvert\braket{n\vert\psi}\rvert^2.
\end{equation}
This is the standard discrete version of the Born rule.\supercite{10.1007/BF01397477}
\subsection{Spectral Decomposition in Infinite Dimensions}\label{sec:spectral_decompositions_in_infinite_dimensions}
In infinite-dimensional Hilbert spaces, self-adjoint operators may have both discrete and continuous parts of the spectrum. The finite sum in the previous subsection is replaced by a combination of sums and integrals. Heuristically, one can write the spectral decomposition of a self-adjoint operator $A$ as\supercite{1021172445}
\begin{equation}\label{eq:heuristic_spectral_decomposition}
A=\sum_n\lambda_n\ket{n}\bra{n}+\int_{\sigma_{\symrm{c}}(A)} \lambda\,\ket{\lambda}\bra{\lambda}\,\symrm{d}\lambda,
\end{equation}
where $\lambda_n$ are discrete eigenvalues with normalizable eigenstates $\ket{n}$ (bound states), $\sigma_{\symrm{c}}(A)$ is the continuous part of the spectrum. The generalized kets $\ket{\lambda}$ are not elements of $\mathcal{H}$ (they are not normalizable), but they can be treated formally in Dirac’s notation. More rigorously, the spectral decomposition is described using a projection-valued measure, but the heuristic form in Eq.~\eqref{eq:heuristic_spectral_decomposition} is sufficient for most physical applications. The probability of obtaining a measurement outcome in an interval $\Delta$ is
\begin{equation}
\mathbb P_A^\psi(\Delta)=\sum_{\lambda_n\in\Delta}\lvert\braket{n\vert\psi}\rvert^2+\int_{\Delta\cap\sigma_{\symrm{c}}(A)}\lvert\braket{\lambda\vert\psi}\rvert^2\,\symrm{d}\lambda,
\end{equation}
where the discrete sum runs over eigenvalues in $\Delta$ and the integral covers the continuous spectrum in $\Delta$.
\subsection{Functions of Observables and Time Evolution}\label{sec:functions_of_observables_and_time_evolution}
Once the spectral decomposition of a self-adjoint operator $A$ is known, one can define functions of $A$ by acting on its spectral components.

In the purely discrete case,
\begin{equation}
A=\sum_n\lambda_n\ket{n}\bra{n},\qquad\ket{\psi}=\sum_n c_n\ket{n},
\end{equation}
a function $f(A)$ is defined by
\begin{equation}
f(A):=\sum_n f(\lambda_n)\ket{n}\bra{n},
\end{equation}
so that
\begin{equation}
f(A)\ket{n}=f(\lambda_n)\ket{n}.
\end{equation}
In the general (discrete + continuous) case, the same idea gives
\begin{equation}
f(A)=\sum_n f(\lambda_n)\ket{n}\bra{n}+\int_{\sigma_{\symrm{c}}(A)}f(\lambda)\ket{\lambda}\bra{\lambda}\,\symrm{d}\lambda,
\end{equation}
which is again a heuristic form. The most important example in quantum mechanics is the time-evolution operator generated by a Hamiltonian $H$. For a time-independent Hamiltonian, the unitary time-evolution operator is
\begin{equation}
U(t)=\e^{-\i tH},
\end{equation}
which, in spectral form, reads
\begin{equation}
U(t)=\sum_n\e^{-\i tE_n} \ket{n}\bra{n}+\int_{\sigma_{\symrm{c}}(H)}\e^{-\i tE}\ket{E}\bra{E}\,\symrm{d}E,
\end{equation}
where $E_n$ and $E$ denote discrete and continuous energy values, respectively. This expression shows explicitly how each energy component of the state picks up a phase factor $\e^{-\i tE}$, while the norm
\begin{equation}
\lvert U(t)\psi\rvert=\lvert\psi\rvert
\end{equation}
is preserved, expressing conservation of total probability under \acrfull{tdse} time evolution.
